\section{Discussion}\label{sec:discussion}

\paragraph{Implications for practice} Three recommendations follow for operators of behind-the-meter storage under demand charges. First, choose and tune a forecaster with an asymmetric score that penalises under-forecasts more than over-forecasts, such as the pinball loss of an upper quantile, not with the MAE. An upper-quantile point forecast is the optimal forecast under such a loss~\cite{gneiting2011point}, it needs no special training method, and here it was the most valuable forecaster although it was the least accurate. Second, never execute a cost-minimising day-ahead schedule open loop: the schedule charges up to the month's peak by design, and any under-forecast at those hours sets a new peak. A guard that refuses to charge into the peak and discharges against it turned every forecaster from a loss into a saving and mattered more than the choice of forecaster. Third, check the bill, not the error, on the operator's own buildings before deploying a forecaster; the attainable saving is a fraction of a per cent of the bill, so small systematic errors decide its sign. Under a pure time-of-use tariff, by contrast, a seasonal naive forecast was as good as any.

\paragraph{Implications for research} The results support the motivation of decision-focused learning~\cite{donti2017task,elmachtoub2022spo,mandi2024dfl} and qualify it. MAE is a poor proxy for value in this decision, but most of the gap is closed by a fixed asymmetric score, without training on the decision. Studies of decision-focused forecasters for storage~\cite{kim2025neural,depoortere2026dfl} should therefore include two baselines: a quantile forecaster, since part of their gain may come from the bias that training on decision cost induces; and a guarded controller, since a gain measured with open-loop execution may disappear once the controller corrects the schedule in real time. Benchmarks of building load forecasting~\cite{emami2023buildingsbench,miller2020gepiii} could report an asymmetric or decision-based score alongside their error metrics, as control benchmarks already score controllers by cost~\cite{vazquez2019citylearn}.

\paragraph{Limitations} The study uses one data set, with buildings from \res{keptSites} sites in 2016--2017, most of them on university campuses, and hourly data, so demand is billed on the hourly rather than the 15-minute peak. The tariff and the battery are stylised and the same for every building, and battery degradation and capital cost are ignored; the savings in dollars indicate magnitudes only. The schedule is a daily linear program that sees the month-to-date peak but not the rest of the month, and the two controllers are simple; a model predictive controller that re-optimises during the day~\cite{mary2025rmpc} would likely narrow the gap between forecasters further. Temperature is observed rather than forecast, the models are not refitted during the test year, and the synthetic errors are deliberately simple. The analysis plan was fixed in the author's own repository rather than with an independent registry, and three amendments were made, the last after the results and a review were seen; the registered test gave the same answer on the registered and the corrected simulation, and the conclusions about the sign of the error rest on secondary and exploratory analyses that a confirmatory study should test on new buildings.
